The theoretical audit supports the intended ELBO interpretation after correcting sign and indexing discrepancies; it does not establish finite-model optimality or rollout stability. Claim 1 receives limited support from eight 80-frame videos with modest early-to-late temporal-change differences and no aggregate collapse. Released-checkpoint maze evaluation shows generally improved reward and reach time with stronger guidance, but cannot establish the paper's claimed advantage over baselines under a matched protocol. The released architecture and planning protocol differ materially from the original experiments, which limits what inference with these checkpoints can reproduce despite a working inference environment.

All eight video conditions, both maze defaults with three model seeds, and both guidance sweeps are complete. Full-sequence sampling is about twice as fast, while temporal-change metrics do not consistently favor either schedule. Visual samples show plausible scene structure together with local artifacts and incomplete maze paths. The final report may examine time-series Table 2 on the original branch, MCG and Diffuser under a matched evaluation protocol, and a small teacher-forcing baseline. None of those tasks was started here.

Retraining the paper's video model is not a credible near-term target on one 4090. Smaller time-series or maze models may be feasible, subject to measured resource needs. MDF-lite depends on demonstration collection meeting the November 19 milestone and training by December 3. Author-contact status and any reply must be confirmed by the team; questions should include the released baselines, robot data, evaluation normalization, and the Appendix A discrepancies documented here.
